% ថ្នាក់ទី ១១ (វិទ្យាសាស្ត្រ) · ជំពូកទី ៨ · មេរៀនទី ៣ · សៀវភៅភាគ ២ ទំព័រ ១៧៨–២០០
\chapter{ភាពអរត្តូកូណាល់ក្នុងលំហ}

\begin{objBox}[វត្ថុបំណង]
\begin{enumerate}[1]
  \item កំណត់និងប្រើប្រាស់លក្ខណៈបន្ទាត់ពីរអរត្តូកូណាល់គ្នាក្នុងការដោះស្រាយលំហាត់
  \item កំណត់និងប្រើប្រាស់លក្ខណៈបន្ទាត់កែងនឹងប្លង់ក្នុងការដោះស្រាយលំហាត់
  \item កំណត់និងប្រើប្រាស់លក្ខណៈប្លង់មេដ្យាទ័រក្នុងការដោះស្រាយលំហាត់
  \item កំណត់និងប្រើប្រាស់លក្ខណៈចំណោលកែងក្នុងការដោះស្រាយលំហាត់ធរណីមាត្រ ។
\end{enumerate}
\end{objBox}

\section{រំលឹករូបមន្តផ្ទៃក្រឡានិងមាឌនៃសូលីតខ្លះៗ}
តាង $V$ ជាមាឌ $A$ ជាផ្ទៃក្រឡាខាង $B$ ជាផ្ទៃក្រឡាបាត $r$ ជាកាំ និង $h$ ជាកម្ពស់ ។
\[
\begin{array}{|l|l|l|}\hline
\text{សូលីត} & \text{មាឌ} & \text{ផ្ទៃក្រឡាខាង}\\ \hline
\text{ប្រលេពីប៉ែតកែង} & V = a\times b\times c & A = 2c(a + b)\\[3pt]
\text{ព្រីសត្រង់} & V = B\times h & A = h(a + b + c + d + e)\\[3pt]
\text{ស៊ីឡាំង} & V = \pi r^2h & A = 2\pi rh\\[6pt]
\text{ស្វ៊ែ} & V = \dfrac{4}{3}\pi R^3 & A = 4\pi R^2\\[6pt]
\text{ពីរ៉ាមីតនិយ័ត} & V = \dfrac{1}{3}Bh & \\[6pt]
\text{កោណ} & V = \dfrac{1}{3}\pi r^2h & A = \pi ra\\ \hline
\end{array}
\]
ក្នុងតារាង ៖ $a, b, c$ ជាវិមាត្រនៃប្រលេពីប៉ែតកែង ($c$ ជាកម្ពស់), $a, b, c, d, e$ ជាជ្រុងនៃបាតព្រីស និង $a$ ជាទ្រនុងខាង (ជនេត្រ) នៃកោណ ។

\begin{example}[លំហាត់គំរូ]
ប្រសិនបើផ្ទៃសរុបនៃគូបមួយស្មើនឹង $150\ \text{m}^2$ ។ ចូររកប្រវែងទ្រនុងនិងមាឌនៃគូប ។
\ans
តាង $a$ ជាប្រវែងទ្រនុងនៃគូប ។ ផ្ទៃក្រឡានៃមុខខាងមួយស្មើនឹង $a^2$ ។ គេបានផ្ទៃក្រឡាសរុបនៃគូបនេះគឺ $6a^2 = 150$, $a^2 = 25$ ឬ $a = \sqrt{25} = 5$ ។ ដូច្នេះ ប្រវែងទ្រនុងនៃគូបគឺ $5$ m ។ មាឌគូបគឺ $V = 5\times 5\times 5 = 125\ \text{m}^3$ ។
\end{example}

\begin{pro}
ព្រីសឆមុខនិយ័តមួយ មានផ្ទៃបាតមួយស្មើ $24\sqrt{3}\ \text{cm}^2$ ។ ប្រសិនបើមុខខាងនីមួយៗជាការេ ។ ចូររកផ្ទៃក្រឡាសរុប និងមាឌនៃព្រីសនោះ ។
\ans
ផ្ទៃបាត $\dfrac{3\sqrt{3}}{2}x^2 = 24\sqrt{3} \Rightarrow x = 4$ cm (ជ្រុងបាត $=$ កម្ពស់ $= 4$ cm) ។ ផ្ទៃក្រឡាសរុប $= 6\times 4^2 + 2\times 24\sqrt{3} = 96 + 48\sqrt{3} = 48(2 + \sqrt{3})\ \text{cm}^2$ ។ មាឌ $= 24\sqrt{3}\times 4 = 96\sqrt{3}\ \text{cm}^3$ ។
\end{pro}

\begin{pro}
គេយកបន្ទះដែកមួយរាងចតុកោណកែង $ABDC$ ($AB = 12$ cm, $AC = 8$ cm) ពត់ធ្វើជាស៊ីឡាំងបរិវត្តតាមជ្រុង $AB$ ។ ចូររកមាឌស៊ីឡាំងនេះ ។
\ans
បរិមាត្របាត $2\pi r = 12 \Rightarrow r = \dfrac{6}{\pi}$ និងកម្ពស់ $h = 8$ ។ មាឌស៊ីឡាំង $= \pi\left(\dfrac{6}{\pi}\right)^2\times 8 = \dfrac{288}{\pi}\ \text{cm}^3$ ។
\end{pro}

\section{បន្ទាត់អរត្តូកូណាល់}
គេមានតុមួយ ($CDFE$ ជាផ្ទៃតុ ហើយ $CC'$, $DD'$ ជាជើងតុ ដែល $AB$ ជាឈើទ្រនាប់ស្របនឹង $CD$) ។ គេសង្កេតឃើញថា $AB \parallel CD$ នោះមុំ $\angle CDD'$ ជាមុំរវាងបន្ទាត់ $AB$ និង $CC'$ ។ ដោយ $\angle CDD'$ ជាមុំកែងនោះគេថាបន្ទាត់ $AB$ អរត្តូកូណាល់នឹង $CC'$ ។ គេកំណត់សរសេរ $AB \perp CC'$ ។

\begin{definition}[និយមន័យ]
បន្ទាត់ពីរនៅក្នុងលំហ \textbf{អរត្តូកូណាល់គ្នា} កាលណាមុំរវាងបន្ទាត់ទាំងពីរជាមុំកែង ។
\end{definition}

\begin{example}[លំហាត់គំរូ]
ក្នុងគូប $ABCDA'B'C'D'$ ៖
\end{example}
\begin{itemize}
  \item បន្ទាត់ដែលស្របនឹង $AA'$ ហើយកាត់តាម $B$ គឺ $BB'$ ។ តែ $BB'$ កែងនឹង $BC$ ដូច្នេះ $AA'$ អរត្តូកូណាល់នឹង $BC$ ។
  \item ដូចគ្នាដែរ $\left.\begin{aligned} BD &\parallel B'D'\\ AC &\perp BD \end{aligned}\right\} \Rightarrow AC$ អរត្តូកូណាល់នឹង $B'D'$ ។
\end{itemize}
\begin{note}[សម្គាល់]
ក្នុងឧទាហរណ៍ខាងលើ គេមានបន្ទាត់ $AC$ និង $BD$ ស្ថិតក្នុងប្លង់ $ABCD$ តែមួយ ហើយ $AC$ កែងនឹង $BD$ ។ គេកំណត់សរសេរ $AC \perp BD$ ។ បន្ទាត់ $AC$ និង $B'D'$ ស្ថិតក្នុងប្លង់ផ្សេងគ្នា ហើយ $AC$ អរត្តូកូណាល់នឹង $B'D'$ ។ គេកំណត់សរសេរអរត្តូកូណាល់នេះ ដោយប្រើសញ្ញា $\perp$ ដែរគឺ $AC \perp B'D'$ ។
\end{note}
\begin{theorem}[វិបាក]
\[
\left.\begin{aligned} a &\perp d\\ a &\parallel b \end{aligned}\right\} \Rightarrow b \perp d
\qquad\qquad
\left.\begin{aligned} a &\perp d\\ b &\perp d\\ a, b, d &\subset P \end{aligned}\right\} \Rightarrow a \parallel b
\]
\end{theorem}

\begin{pro}
គេមានចតុមុខ $ABCD$ ដែល $DA$ អរត្តូកូណាល់ទៅនឹង $BC$ ។ ដៅ $I, J, K$ ជាចំណុចកណ្ដាលរៀងគ្នានៃជ្រុង $AB$, $BD$, $CD$ ។
\begin{enumerate}[a]
  \item បង្ហាញថា $JK$ អរត្តូកូណាល់ទៅនឹង $DA$ ។
  \item បង្ហាញថា $IJ$ អរត្តូកូណាល់ទៅនឹង $BC$ (ហើយ $IJ \perp JK$) ។
\end{enumerate}
\ans
ក. ក្នុងត្រីកោណ $BCD$ ៖ $JK \parallel BC$ ។ ដោយ $BC \perp DA$ នោះ $JK \perp DA$ ។ \quad ខ. ក្នុងត្រីកោណ $ABD$ ៖ $IJ \parallel AD$ ។ ដោយ $AD \perp BC$ នោះ $IJ \perp BC$ ហើយដោយ $JK \parallel BC$ នោះ $IJ \perp JK$ ។
\end{pro}

\section{បន្ទាត់កែងនឹងប្លង់}
ប្រសិនបើគេយកសៀវភៅមួយក្បាលដាក់បញ្ឈរពីលើប្លង់តុដេកមួយ នោះគេឃើញថា សន្លឹកសៀវភៅទាំងអស់ប្រសព្វគ្នាបានបន្ទាត់ $AA'$ ។ គែមនៃសន្លឹកសៀវភៅស្ថិតជាប់នឹងប្លង់តុដូចជា $AB$, $AC$, $\dots$ ។ ដូច្នេះនៅក្នុងប្លង់តុនេះមានបន្ទាត់ជាច្រើន $AB$, $AC$, $\dots$ ដែលកែងនឹងបន្ទាត់ $AA'$ ត្រង់ចំណុច $A$ ។ ម្យ៉ាងទៀតបើគេយកបន្ទាត់ $b$ មួយផ្សេង មកដាក់ដេកលើតុ នោះគេសង្កេតឃើញថា $AA'$ អរត្តូកូណាល់នឹងបន្ទាត់ $b$ ។
\begin{definition}[និយមន័យ]
បន្ទាត់មួយ\textbf{កែងនឹងប្លង់}មួយ កាលណាបន្ទាត់នោះអរត្តូកូណាល់នឹងគ្រប់បន្ទាត់នៅក្នុងប្លង់ ។
\end{definition}
តាមនិយមន័យខាងលើនេះ គេអាចជ្រើសរើសយកតែបន្ទាត់ប្រសព្វគ្នាពីរនៅក្នុងប្លង់នោះ គេបានទ្រឹស្ដីបទដូចខាងក្រោម ៖
\begin{theorem}[ទ្រឹស្ដីបទ ១]
បើបន្ទាត់មួយ អរត្តូកូណាល់នឹងបន្ទាត់ពីរប្រសព្វគ្នានៅក្នុងប្លង់ នោះវាកែងនឹងប្លង់ ។
\end{theorem}
\textbf{សម្រាយបញ្ជាក់ ៖} គេមានបន្ទាត់ $a$ និង $b$ នៅក្នុងប្លង់ $\alpha$ ដែល $a \cap b = \{O\}$ ហើយបន្ទាត់ $d$ កែងនឹងបន្ទាត់ $a$ និង $b$ ត្រង់ $O$ ។ គេត្រូវបង្ហាញថា $d$ កែងនឹងគ្រប់បន្ទាត់ $c$ នៃប្លង់ $\alpha$ ដែលកាត់តាម $O$ ។ គូសបន្ទាត់ $m$ មួយនៅក្នុងប្លង់ $\alpha$ ដែលមិនកាត់តាមចំណុច $O$ តែកាត់បន្ទាត់ $a$, $b$ និង $c$ រៀងគ្នាត្រង់ចំណុច $A$, $B$ និង $C$ ។ នៅលើបន្ទាត់ $d$ ជ្រើសរើសចំណុច $P$ មួយ រួចបន្លាយ $PO$ ឱ្យបាន $OQ$ ដែល $PO = OQ$ ។ គេបានត្រីកោណកែង ៖
\[
\begin{aligned}
&\triangle PAO \cong \triangle QAO \Rightarrow AP = AQ && (1)\\
&\triangle PBO \cong \triangle QBO \Rightarrow BP = BQ && (2)\\
&(1), (2) \Rightarrow \triangle PAB \cong \triangle QAB \Rightarrow \angle PAB\\
&= \angle QAB && (3)
\end{aligned}
\]
(1) និង (3) គេមាន $\triangle PAC \cong \triangle QAC$ នោះ $PC = QC$ ។ គេបាន $\triangle PCQ$ ជាត្រីកោណសមបាតដោយ $O$ ជាចំណុចកណ្ដាល $PQ$ នោះ $OP \perp OC$ ពោលគឺ $d \perp c$ ។

\begin{note}[សម្គាល់]
បើបន្ទាត់ $d$ កែងនឹងបន្ទាត់ $a$ និង $b$ មិននៅត្រង់ចំណុចប្រសព្វ $O$ រវាងបន្ទាត់ទាំងពីរ ក្នុងករណីនេះគេគូសបន្ទាត់ $d'$ មួយកាត់តាម $O$ ដែល $d' \parallel d$ គេបាន បើ $d' \perp c$ នោះ $d \perp c$ ដែរ ។
\end{note}

\begin{example}[លំហាត់គំរូ]
គេឱ្យប្រលេពីប៉ែតកែង $ABCDEFGH$ ($ABCD$ ជាបាត ដែល $AB = a$, $AD = b$, $AE = c$ និង $E, F, G, H$ នៅពីលើ $A, B, C, D$ រៀងគ្នា) ។
\begin{enumerate}[a]
  \item បង្ហាញថាបន្ទាត់ $BC$ អរត្តូកូណាល់នឹងមុខ $ABFE$
  \item បង្ហាញថាមុំ $\angle EBC$ ជាមុំកែង
  \item រកប្រវែង $EC$ ជាអនុគមន៍នឹង $a$, $b$ និង $c$ ។
\end{enumerate}
\ans
\begin{enumerate}[a]
  \item ក្នុងចតុកោណកែង $ABCD$ មាន $BC \perp AB$ ។ ក្នុងចតុកោណកែង $BCGF$ មាន $BC \perp BF$ ។ ដូច្នេះ $BC$ កែងនឹងប្លង់ $ABFE$ មានន័យថា $BC$ អរត្តូកូណាល់នឹងមុខ $ABFE$ ។
  \item $\left.\begin{aligned} BC &\perp ABFE\\ EB &\subset ABFE \end{aligned}\right\} \Rightarrow BC \perp EB$ ។ ដូច្នេះមុំ $\angle EBC$ ជាមុំកែង ។
  \item តាមសំណួរ ខ ៖ $\triangle EBC$ ជាត្រីកោណកែងត្រង់ $B$ ។ ដូច្នេះ $EC^2 = EB^2 + BC^2 = EB^2 + b^2$ ។ ដោយ $\triangle ABE$ ជាត្រីកោណកែងត្រង់ $A$ នោះ $EB^2 = AE^2 + AB^2 = c^2 + a^2$ ។ គេបាន $EC^2 = a^2 + b^2 + c^2$ ។ ដូច្នេះ $EC = \sqrt{a^2 + b^2 + c^2}$ ។
\end{enumerate}
\end{example}

\begin{pro}
គេឱ្យបន្ទាត់ $AK$ កែងទៅនឹងប្លង់ $Q$ ត្រង់ $K$ និងបន្ទាត់ $d$ មួយឋិតក្នុងប្លង់ $Q$ ប៉ុន្តែមិនកាត់តាម $K$ ។ ដៅចំណុច $C$ មួយនៅលើបន្ទាត់ $d$ ដែលបន្ទាត់ $KC$ កែងទៅនឹងបន្ទាត់ $d$ ។ យក $B$ ជាចំណុចមួយឋិតនៅលើបន្ទាត់ $d$ ។
\begin{enumerate}[a]
  \item បង្ហាញថា $d$ កែងទៅនឹងប្លង់ $AKC$
  \item បង្ហាញថា $\angle ACB$ ជាមុំកែង
  \item បង្ហាញថា $AB^2 = AK^2 + KC^2 + BC^2$ ។
\end{enumerate}
\ans
ក. $d \perp KC$ និង $d \perp AK$ (ព្រោះ $AK \perp Q$) ដែល $KC$ និង $AK$ ប្រសព្វគ្នា ដូច្នេះ $d \perp (AKC)$ ។ \quad ខ. $AC \subset (AKC)$ នាំឱ្យ $d \perp AC$ គឺ $\angle ACB = 90^\circ$ ។ \quad គ. $AB^2 = AC^2 + CB^2$ និង $AC^2 = AK^2 + KC^2$ ដូច្នេះ $AB^2 = AK^2 + KC^2 + BC^2$ ។
\end{pro}

\begin{theorem}[ទ្រឹស្ដីបទទី ២]
\begin{enumerate}[a]
  \item $\left.\begin{aligned} a &\parallel b\\ a &\perp P \end{aligned}\right\} \Rightarrow b \perp P$
  \item $\left.\begin{aligned} a &\perp P\\ b &\perp P \end{aligned}\right\} \Rightarrow a \parallel b$
  \item $\left.\begin{aligned} a &\perp P\\ P &\parallel Q \end{aligned}\right\} \Rightarrow a \perp Q$
\end{enumerate}
\end{theorem}
\textbf{សម្រាយបញ្ជាក់ ៖}
\begin{enumerate}[a]
  \item គេឱ្យបន្ទាត់ $a$ មួយអរត្តូកូណាល់ និងបន្ទាត់ប្រសព្វពីរ $c$ និង $d$ នៃប្លង់ $P$ ត្រង់ $O$ ។ បើគេមានបន្ទាត់ $b \parallel a$ នោះគេបានអរត្តូកូណាល់នឹង $c$ និង $d$ ។ ដោយ $a$ អរត្តូកូណាល់នឹងគ្រប់បន្ទាត់ក្នុងប្លង់និង $a \parallel b$ ដូច្នេះ គេបាន $b$ កែងនឹងប្លង់ $P$ ។
  \item ឧបមាថាបន្ទាត់ $b \nparallel a$ ដូច្នេះគេអាចរកបាន $b'$ ដែល $b' \parallel a$ ហើយ $b'$ កាត់តាម $B$ ។ ដោយ $a \perp P$ និង $b' \parallel a$ នោះ $b' \perp P$ ។ តាមទ្រឹស្ដីបទ (ក) តាម $B$ មានបន្ទាត់ $b$ និង $b'$ ដែលកែងនឹងប្លង់ $P$ តែមួយជាករណីមិនអាចមាន ។ ដូច្នេះ គេបាន $b \parallel a$ ។
  \item ដោយបន្ទាត់ $a$ កែងនឹង $P$ នោះ $a$ អរត្តូកូណាល់នឹងគ្រប់បន្ទាត់ក្នុងប្លង់ $P$ ។ ជ្រើសរើសបន្ទាត់ពីរ $b$ និង $c$ ប្រសព្វគ្នាត្រង់ $O$ ដែលនៅក្នុងប្លង់ $P$ ។ ដូច្នេះ $a$ កែងនឹង $b$ និង $c$ ត្រង់ $O$ ។ ម្យ៉ាងទៀត $P \parallel Q$ នាំឱ្យមានបន្ទាត់ $b' \parallel b$, $c' \parallel c$ ដែល $b' \subset Q$, $c' \subset Q$ និង $b' \cap c' = \{O'\}$ គេបាន $a$ កែងនឹង $b'$ និង $c'$ ត្រង់ $O'$ ។ តាមទ្រឹស្ដីបទ ១ គេបាន $a \perp Q$ ។
\end{enumerate}

\begin{example}[លំហាត់គំរូ]
គេឱ្យ $4$ ចំណុច $S, A, B$ និង $C$ មិនស្ថិតក្នុងប្លង់តែមួយ ដែល $SA$ កែងនឹងប្លង់ $ABC$ ។ $BM$, $BN$ និង $SO$ កែងរៀងគ្នានឹង $SC$, $AC$ និង $BC$ ។ $E$, $F$ ជាចំណុចប្រសព្វរៀងគ្នានៃ $AO$ និង $BN$, $SO$ និង $BM$ ។ បង្ហាញថា ៖
\begin{enumerate}[a]
  \item $AO$ កែងនឹង $BC$
  \item $SC$ កែងនឹងប្លង់ $BMN$
  \item $EF$ កែងនឹងប្លង់ $SBC$ ។
\end{enumerate}
\ans
\begin{enumerate}[a]
  \item $SA \perp ABC$ និង $BC \subset ABC$ នាំឱ្យ $SA \perp BC$ ។ ដោយ $SO \perp BC$ ផង នោះ $BC \perp (SOA)$ ។ ហើយ $AO \subset (SOA)$ ដូច្នេះ $BC \perp AO$ ។
  \item $SA \perp ABC$ និង $BN \subset ABC$ នាំឱ្យ $SA \perp BN$ ។ ដោយ $BN \perp AC$ ផង នោះ $BN \perp (SAC)$ ដូច្នេះ $BN \perp SC$ ។ ម្យ៉ាងទៀត $BM \perp SC$ ។ ហេតុនេះ $SC \perp (BMN)$ ។
  \item $E \in AO$, $E \in BN$ និង $F \in SO$, $F \in BM$ នាំឱ្យ $EF \subset (SOA)$ និង $EF \subset (BMN)$ ។ ដោយ $SC \perp (BMN)$ នោះ $SC \perp EF$ \quad (1) ។ ដោយ $BC \perp (SOA)$ នោះ $BC \perp EF$ \quad (2) ។ តាម (1) និង (2) ៖ $EF$ កែងនឹងបន្ទាត់ប្រសព្វពីរ $SC$ និង $BC$ នៃប្លង់ $SBC$ ដូច្នេះ $EF \perp (SBC)$ ។
\end{enumerate}
\end{example}

\begin{pro}
គេមានចតុមុខនិយ័ត $OABC$ ។ $H$ ជាអរត្តូសង់នៃ $\triangle ABC$ ។ $K, M, N$ និង $L$ ជាចំណុចកណ្ដាលរៀងគ្នានៃ $OB$, $OC$, $HB$ និង $HC$ ។
\begin{enumerate}[a]
  \item បង្ហាញថា $OH \perp ABC$
  \item បន្ទាត់ $a$ ជាប្រសព្វនៃ $AKN$ និង $AML$ ។ បង្ហាញថា $a$ កែងនឹង $ABC$ ។
\end{enumerate}
\ans
ក. ក្នុងចតុមុខនិយ័ត ទ្រនុងឈម $OA \perp BC$ ។ ដោយ $AH \perp BC$ នោះ $BC \perp (OAH)$ នាំឱ្យ $BC \perp OH$ ។ ដូចគ្នា $AC \perp OH$ ។ ដូច្នេះ $OH \perp ABC$ ។ \quad ខ. $KN \parallel OH$ (បាតមធ្យមនៃ $\triangle OBH$) និង $ML \parallel OH$ (បាតមធ្យមនៃ $\triangle OCH$) ។ ដូច្នេះ $OH$ ស្របនឹងប្លង់ $AKN$ និង $AML$ ហើយ $a \parallel OH$ ។ តាមទ្រឹស្ដីបទទី ២ (ក) $a \perp ABC$ ។
\end{pro}

\begin{theorem}[ទ្រឹស្ដីបទបីបន្ទាត់កែង]
បើបន្ទាត់ $AB$ កែងនឹងប្លង់ $P$ ត្រង់ $B$ ហើយ $BO$ កែងនឹងបន្ទាត់ $a$ នៃប្លង់ $P$ ត្រង់ $O$ នោះបន្ទាត់ $AO$ កែងនឹងបន្ទាត់ $a$ ៖
\[
\left.\begin{aligned} AB &\perp P\\ BO &\perp a\\ a &\subset P \end{aligned}\right\} \Rightarrow AO \perp a
\]
\end{theorem}
\textbf{សម្រាយបញ្ជាក់ ៖} $AB \perp P$, $a \subset P$ នាំឱ្យ $AB \perp a$ ។ $AB \perp a$, $BO \perp a$ នាំឱ្យ $a \perp (AOB)$ ។ ដោយ $AO \subset (AOB)$ នោះ $AO \perp a$ ។

\begin{example}[លំហាត់គំរូ]
គេមានចតុមុខនិយ័ត $SABC$ មួយដែលមានកម្ពស់ $SH$ ។ យក $K$ ជាចំណុចកណ្ដាលនៃជ្រុង $BC$ ។ បង្ហាញថា $SK$ កែងនឹង $BC$ ។
\ans
តាមទ្រឹស្ដីបទបីបន្ទាត់កែង ៖ $SH \perp ABC$, $AK \perp BC$ ($H \in AK$) និង $BC \subset ABC$ នោះ $SK \perp BC$ ។
\end{example}

\begin{theorem}[ទ្រឹស្ដីបទច្រាសនៃបីបន្ទាត់កែង]
បើបន្ទាត់ $AB$ កែងនឹងប្លង់ $P$ ត្រង់ $B$ ហើយ $AO$ កែងនឹងបន្ទាត់ $a$ នៃប្លង់ $P$ ត្រង់ $O$ នោះបន្ទាត់ $BO$ កែងនឹងបន្ទាត់ $a$ ៖
\[
\left.\begin{aligned} AB &\perp P\\ AO &\perp a\\ a &\subset P \end{aligned}\right\} \Rightarrow BO \perp a
\]
\end{theorem}
\textbf{សម្រាយបញ្ជាក់ ៖} $AB \perp P$, $a \subset P$ នាំឱ្យ $AB \perp a$ ។ $AB \perp a$, $AO \perp a$ នាំឱ្យ $a \perp (AOB)$ ។ ដោយ $BO \subset (AOB)$ នោះ $BO \perp a$ ។
\begin{note}[សម្គាល់]
បន្ទាត់ $AB$, $AO$ និង $a$ ដែលកែងរវាងគ្នាពីរៗ ហៅថា \textbf{បីបន្ទាត់កែង} ។
\end{note}

\begin{pro}
គេមានពីរ៉ាមីតចតុមុខនិយ័ត $SABCD$ ដែលមានបាត $ABCD$ ជាការេ ។ យក $M$ ជាចំណុចកណ្ដាលនៃជ្រុង $BC$ ។
\begin{enumerate}[a]
  \item បង្ហាញថា $SM$ កែងនឹង $BC$
  \item ប្រសិនបើ $SA = 20$ cm និង $AB = 5$ cm គណនាផ្ទៃក្រឡាខាងនៃពីរ៉ាមីត ។
\end{enumerate}
\ans
ក. $SB = SC$ ដូច្នេះ $\triangle SBC$ ជាត្រីកោណសមបាត ហើយមេដ្យាន $SM$ ក៏ជាកម្ពស់ ៖ $SM \perp BC$ (ឬប្រើទ្រឹស្ដីបទបីបន្ទាត់កែង) ។ \quad ខ. $SM = \sqrt{20^2 - 2.5^2} = \dfrac{15\sqrt{7}}{2}$ ។ ផ្ទៃក្រឡាខាង $= 4\times\dfrac{1}{2}\times 5\times\dfrac{15\sqrt{7}}{2} = 75\sqrt{7}\ \text{cm}^2$ ។
\end{pro}

\section{ប្លង់មេដ្យាទ័រ}
គេមានពីរចំណុចផ្សេងគ្នា $A$ និង $B$ ហើយ $O$ ជាចំណុចកណ្ដាលនៃអង្កត់ $AB$ ។ បើមានប្លង់ $P$ មួយកែងនឹង $AB$ ត្រង់ចំណុចកណ្ដាល $O$ នោះប្លង់ $P$ នេះហៅថា \textbf{ប្លង់មេដ្យាទ័រ}នៃអង្កត់ $AB$ ។
\begin{definition}[និយមន័យ]
ប្លង់មេដ្យាទ័រនៃអង្កត់មួយគឺជាប្លង់ដែលកែងនឹងអង្កត់នោះត្រង់ចំណុចកណ្ដាល ។
\end{definition}
ឧបមាថា ប្លង់ $P$ ជាប្លង់មេដ្យាទ័រនៃអង្កត់ $AB$ កាត់ត្រង់ $O$ ហើយ $C$ ជាចំណុចមួយក្នុងប្លង់ $P$ ។ បើគេមានចំណុច $C'$ មួយស្ថិតនៅខាង $A$ នោះគេបាន $AC' < BC'$ ។ ម្យ៉ាងទៀតបើគេមានចំណុច $C'$ មួយស្ថិតនៅខាង $B$ នោះគេបាន $BC' < AC'$ ។ តាមលទ្ធផលនេះបង្ហាញថា $AC'$ មិនស្មើនឹង $BC'$ ។ ដូច្នេះ គេបាន $CA = CB$ ដែល $C \in P$ ។
\begin{theorem}[លក្ខណៈ]
គេមានចំណុច $A$, $B$ ស្ថិតនៅសងខាងប្លង់ $P$ ហើយបីចំណុចផ្សេងគ្នា $C$, $D$, $E$ (មិនរត់ត្រង់គ្នា) ស្ថិតនៅក្នុងប្លង់ $P$ ដែល ៖
\[
\left.\begin{aligned} AC &= BC\\ AD &= BD\\ AE &= BE \end{aligned}\right\} \Leftrightarrow P \text{ ជាប្លង់មេដ្យាទ័រនៃអង្កត់ } AB
\]
\end{theorem}

\begin{example}[លំហាត់គំរូ]
ក្នុងគូប $ABCDA'B'C'D'$ ។ ចូរកំណត់ប្លង់មេដ្យាទ័រនៃអង្កត់ $[AC]$ ។
\ans
ក្នុងគូប $ABCDA'B'C'D'$ គេបាន $AB = CB$, $AB' = CB'$, $AD = CD$ ។ ដោយចំណុច $B$, $B'$ និង $D$ មិនរត់ត្រង់គ្នា ដូច្នេះ ប្លង់មេដ្យាទ័រនៃអង្កត់ $AC$ គឺប្លង់ $BB'D$ ។
\end{example}

\begin{pro}
ដោយប្រើគូបនៃឧទាហរណ៍ខាងលើ ហើយយក $I$, $J$, $K$ ជាចំណុចកណ្ដាលរៀងគ្នានៃទ្រនុង $BC$, $CD$, $A'B'$ ។ កំណត់ប្លង់មេដ្យាទ័រនៃអង្កត់ $AC'$ ។
\ans
ប្លង់មេដ្យាទ័រនៃអង្កត់ $AC'$ គឺ $IJK$ (ព្រោះ $IA = IC'$, $JA = JC'$, $KA = KC'$) ។
\end{pro}

\section{ប្លង់កែង}
\begin{definition}[និយមន័យ]
គេថាប្លង់ $P$ \textbf{កែងនឹងប្លង់} $Q$ កាលណាប្លង់ $P$ មានបន្ទាត់មួយកែងនឹងប្លង់ $Q$ ឬប្លង់ $Q$ មានបន្ទាត់មួយកែងនឹងប្លង់ $P$ ៖
\[
\left.\begin{aligned} a &\subset P\\ a &\perp Q \end{aligned}\right\} \Rightarrow P \perp Q
\]
\end{definition}
\begin{theorem}[លក្ខណៈ]
\[
\text{(ក)}\ \left.\begin{aligned} d &\perp P\\ d &\parallel Q \end{aligned}\right\} \Rightarrow P \perp Q
\qquad\qquad
\text{(ខ)}\ \left.\begin{aligned} P &\perp R\\ R &\parallel Q \end{aligned}\right\} \Rightarrow P \perp Q
\]
\end{theorem}
\textbf{សម្រាយបញ្ជាក់លក្ខណៈ (ក) ៖} គេមាន $d \perp P$ និង $d \parallel Q$ ។ នៅក្នុងប្លង់ $Q$ មានបន្ទាត់ $b$ ដែលស្របនឹងបន្ទាត់ $d$ ។ តាមទ្រឹស្ដីបទ ២ (ក) នៃបន្ទាត់កែងនឹងប្លង់ ៖ $d \parallel b$, $d \perp P$ នាំឱ្យ $b \perp P$ ។ ដូច្នេះ $b \subset Q$, $b \perp P$ នាំឱ្យ $P \perp Q$ ។

\begin{pro}
គេឱ្យ $4$ ចំណុច $A$, $B$, $C$ និង $D$ មិនឋិតក្នុងប្លង់តែមួយដែលត្រីកោណ $ABC$ និង $DBC$ ជាត្រីកោណសមបាតមានបាតរួម $BC$ ។ គូស $BO$ កែងនឹង $AD$ ត្រង់ $O$ ។ យក $M$, $N$, $P$ ជាចំណុចកណ្ដាលរៀងគ្នា $AB$, $BD$ និង $BC$ ។
\begin{enumerate}[a]
  \item បង្ហាញថា $BC$ កែងនឹង $APD$ រួច $AD$ កែងនឹង $BCO$ ។
  \item បង្ហាញថា $CMN$ កែងនឹង $BCO$ រួច $ABC$ កែងនឹង $APD$ ។
\end{enumerate}
\ans
ក. $AP \perp BC$ និង $DP \perp BC$ (មេដ្យាននៃត្រីកោណសមបាត) ដូច្នេះ $BC \perp (APD)$ ។ ហេតុនេះ $BC \perp AD$ ហើយដោយ $BO \perp AD$ នោះ $AD \perp (BCO)$ ។ \quad ខ. $MN \parallel AD$ (បាតមធ្យមនៃ $\triangle ABD$) ដូច្នេះ $MN \perp (BCO)$ ហើយប្លង់ $CMN$ ដែលមាន $MN$ កែងនឹង $BCO$ ។ ម្យ៉ាងទៀត $BC \subset ABC$ និង $BC \perp (APD)$ ដូច្នេះ $ABC \perp APD$ ។
\end{pro}

\section{មុំដែលបង្កើតឡើងដោយប្លង់ពីរ}
គេជ្រើសរើសចំណុច $O$ មួយឋិតនៅលើបន្ទាត់ $d$ ដែលជាប្រសព្វរវាងប្លង់ $P$ និង $Q$ រួចគូសបន្ទាត់ $OA$ និង $OB$ កែងទៅនឹង $d$ ដែលចំណុច $A$ និង $B$ នៅក្នុងប្លង់ $P$ និង $Q$ រៀងគ្នា ។ គេបានមុំ $\angle AOB$ ជាមុំដែលបង្កើតឡើងដោយប្លង់ $P$ និង $Q$ ។

\begin{example}[លំហាត់គំរូ]
គេមានគូប $ABCDEFGH$ មួយ ($ABCD$ ជាបាត ហើយ $E, F, G, H$ នៅពីលើ $A, B, C, D$ រៀងគ្នា) ។ កំណត់មុំដែលបង្កើតឡើងដោយប្លង់ $AFGD$ និង $ABCD$ ។
\ans
ប្លង់ $ABCD$ និង $AFGD$ ប្រសព្វគ្នាបានបន្ទាត់ $AD$ ។ គេបាន $BA$ កែងនឹង $AD$ ត្រង់ $A$ និង $AF$ កែងនឹង $AD$ ត្រង់ $A$ ។ ដូច្នេះ គេបានមុំដែលបង្កើតឡើងដោយប្លង់ $AFGD$ និង $ABCD$ គឺជាមុំ $\angle BAF = 45^\circ$ ។
\end{example}

\begin{pro}
គេមានចតុមុខនិយ័ត $OABC$ មួយ និង $M$ ជាចំណុចកណ្ដាល $BC$ ដែល $\angle OAM = \alpha$ ។ កំណត់មុំដែលបង្កើតឡើងដោយប្លង់ $ABC$ និង $OBC$ ។
\ans
$AM \perp BC$ និង $OM \perp BC$ ដូច្នេះមុំរវាងប្លង់ទាំងពីរគឺ $\angle AMO$ ។ ដោយ $AM = OM$ ត្រីកោណ $OAM$ សមបាត ដូច្នេះ $\angle AMO = 180^\circ - 2\alpha$ ។
\end{pro}

\section{ចំណោលកែង}
តាមចំណុច $M$ មួយមិនឋិតនៅក្នុងប្លង់ $P$ គេអាចគូសបានបន្ទាត់តែមួយគត់កាត់តាមចំណុច $M$ ហើយកែងនឹងប្លង់ $P$ ។ តាងចំណុច $M'$ ជាប្រសព្វរវាងបន្ទាត់កែង $MM'$ នឹងប្លង់ $P$ ដូច្នេះ គេហៅថាចំណុច $M'$ ជា\textbf{ចំណោលកែង}នៃ $M$ មកលើប្លង់ $P$ ។
\begin{note}[សម្គាល់]
\begin{itemize}
  \item បើ $M \in P$ នោះ $M'$ ត្រួតលើ $M$ ។
  \item បើ $A'$ ជាចំណោលកែងនៃ $A$ មកលើប្លង់ $P$ នោះ $AA' \perp P$ ។
\end{itemize}
\end{note}
គេមានប្រលេពីប៉ែតកែង $ABCDA'B'C'D'$ ៖
\begin{itemize}
  \item ចំណោលកែងនៃចំណុច $B'$ មកលើប្លង់ $ABC$ គឺជាចំណុច $B$ ។
  \item ចំណោលកែងនៃចំណុច $D'$ មកលើប្លង់ $AA'B'$ គឺជាចំណុច $A'$ ។
  \item ចំណោលកែងនៃចំណុច $O$ ដែលជាចំណុចប្រសព្វរវាង $AC$ និង $BD$ មកលើប្លង់ $A'B'C'D'$ គឺជាចំណុច $O'$ ដែលជាចំណុចប្រសព្វរវាង $A'C'$ និង $B'D'$ ។
\end{itemize}
\begin{definition}[និយមន័យ]
ចំណោលកែងនៃចំណុចមួយមកលើប្លង់មួយ គឺជាចំណុចប្រសព្វរវាងប្លង់និងបន្ទាត់កែងដែលគូសចេញពីចំណុចនោះ ។
\end{definition}
\begin{theorem}[លក្ខណៈ]
\begin{enumerate}[a]
  \item ចំណោលកែងនៃចំណុចរត់ត្រង់គ្នានៅក្រៅប្លង់ $P$ មកលើប្លង់ $P$ គឺជាចំណុចរត់ត្រង់គ្នា ឬត្រួតស៊ីគ្នានៅក្នុងប្លង់ $P$ ។
  \item បើ $A'B'$ ជាចំណោលកែងនៃ $AB$ មកលើប្លង់ $P$ ។ គេបាន $A'B' = AB\cos\alpha$ ដែល $\alpha$ ជាមុំផ្គុំរវាងបន្ទាត់ $AB$ និងប្លង់ $P$ ។
\end{enumerate}
\end{theorem}
\textbf{សម្រាយបញ្ជាក់ ៖}
\begin{enumerate}[a]
  \item គេដឹងមកហើយថាបីចំណុច $A$, $B$, $C$ រត់ត្រង់គ្នា នោះវាឋិតនៅលើបន្ទាត់តែមួយ ។ បើបន្ទាត់ $AB$ កែងនឹងប្លង់ $P$ ត្រង់ចំណុច $A'$ នោះចំណោលកែងនៃ $A$, $B$, $C$ ជាចំណុច $A'$ តែមួយ ។ បើបន្ទាត់ $AB$ មិនកែងនឹងប្លង់ $P$ ៖ តាង $A'$, $B'$, $C'$ ជាចំណោលកែងនៃចំណុច $A$, $B$, $C$ មកលើប្លង់ $P$ រៀងគ្នា ។ គេបានបន្ទាត់ $AA' \parallel BB'$ ហើយកំណត់បានប្លង់ $\alpha$ មួយដែល $P \cap \alpha = A'B'$ ។ ដោយ $CC' \perp P$ និង $AA' \perp P$ នោះ $CC' \parallel AA'$ ។ ម្យ៉ាងទៀត $C \in AB \subset \alpha$ នោះ $C \in \alpha$ ហើយ $CC' \subset \alpha$ មានន័យថា $C' \in \alpha$ ។ គេបាន $C' \in \alpha$ និង $C' \in P$ នាំឱ្យ $C' \in A'B'$ ។ ដូច្នេះ $A'$, $B'$, $C'$ រត់ត្រង់គ្នា ។
  \item គេបាន $AA' \perp A'B'$ ត្រង់ $A'$ និង $BB' \perp A'B'$ ត្រង់ $B'$ ។ យក $AH \perp BB'$ ត្រង់ $H$ នោះ $AH \parallel A'B'$ ហើយ $AH = A'B'$ ។ ដូច្នេះ $\angle HAB = \alpha$ ដែលជាមុំផ្គុំរវាងបន្ទាត់ $AB$ នឹងប្លង់ $P$ ។ ក្នុងត្រីកោណកែង $AHB$ គេបាន $AH = AB\cos\alpha$ ។ ដូច្នេះ $A'B' = AB\cos\alpha$ ។
\end{enumerate}

\begin{pro}
គេមានត្រីកោណ $ABC$ មួយដែល $BC$ ឋិតក្នុងប្លង់ $P$ និងកំពូល $A$ នៅក្រៅប្លង់ $P$ ។ គេធ្វើចំណោលកែងនៃចំណុច $A$ មកលើប្លង់ $P$ បានចំណុច $A'$ ។ ប្រសិនបើមុំផ្គុំរវាងត្រីកោណ $ABC$ នឹងប្លង់ $P$ ស្មើនឹង $\alpha$ ហើយតាង $S$ និង $S'$ ជាផ្ទៃក្រឡានៃត្រីកោណ $ABC$ និង $A'BC$ រៀងគ្នា ។ បង្ហាញថា $S' = S\cos\alpha$ ។
\ans
តាង $AM$ ជាកម្ពស់នៃ $\triangle ABC$ ($M \in BC$) ។ តាមទ្រឹស្ដីបទច្រាសនៃបីបន្ទាត់កែង $A'M \perp BC$ ហើយ $\angle AMA' = \alpha$ ។ $S = \dfrac{1}{2}AM\times BC$, $S' = \dfrac{1}{2}A'M\times BC$ និង $A'M = AM\cos\alpha$ ។ ដូច្នេះ $S' = S\cos\alpha$ ។
\end{pro}

\section*{មេរៀនសង្ខេប ជំពូកទី ៨}
\begin{note}[មេរៀនសង្ខេប]
\begin{enumerate}[1]
  \item ដើម្បីតាងរូបក្នុងលំហ គេត្រូវស្គាល់ច្បាស់អំពីទ្រង់ទ្រាយដើម ។
  \item បន្ទាត់ស្របគ្នាក្នុងលំហ តាងដោយបន្ទាត់ស្របគ្នាក្នុងប្លង់ ។
  \item អង្កត់ពីរដែលមានប្រវែងស្មើគ្នានៅក្នុងលំហអាចមិនចាំបាច់តាងដោយអង្កត់ពីរដែលមានប្រវែងស្មើគ្នានៅក្នុងប្លង់ទេ ។
  \item តាមពីរចំណុចផ្សេងគ្នា គេអាចគូសបន្ទាត់ត្រង់បានតែមួយគត់ ។
  \item តាមបីចំណុចមិនស្ថិតលើបន្ទាត់តែមួយ គេអាចគូសបានប្លង់តែមួយគត់ដែលកាត់តាមបីចំណុចនោះ ។
  \item បន្ទាត់មួយ និងចំណុចមួយមិនស្ថិតលើបន្ទាត់តែមួយ កំណត់បានប្លង់មួយគត់ ។
  \item បន្ទាត់ពីរប្រសព្វគ្នា កំណត់បានប្លង់មួយ ។
  \item បន្ទាត់ពីរស្របគ្នា កំណត់បានប្លង់មួយ ។
  \item បើបន្ទាត់មួយ និងប្លង់មួយមានចំណុចរួមពីរផ្សេងគ្នា នោះបន្ទាត់ទាំងមូលស្ថិតក្នុងប្លង់នោះ ។
  \item បើប្លង់ពីរមានចំណុចរួមមួយ វាប្រសព្វគ្នាតាមបន្ទាត់មួយដែលកាត់តាមចំណុចរួមនោះ ។
  \item \textbf{បន្ទាត់អរត្តូកូណាល់ ៖} បន្ទាត់ពីរនៅក្នុងលំហអរត្តូកូណាល់ឬកែងគ្នា កាលណាមុំរវាងបន្ទាត់ទាំងពីរជាមុំកែង ។
  \item \textbf{បន្ទាត់កែងនឹងប្លង់ ៖} បន្ទាត់មួយកែងនឹងប្លង់មួយ កាលណាបន្ទាត់នោះអរត្តូកូណាល់នឹងគ្រប់បន្ទាត់ក្នុងប្លង់ $P$ ។ បើបន្ទាត់មួយ អរត្តូកូណាល់ឬកែងនឹងបន្ទាត់ពីរប្រសព្វគ្នានៅក្នុងប្លង់ នោះវាកែងនឹងប្លង់ ។
  \item \textbf{ទ្រឹស្ដីបទបីបន្ទាត់កែង ៖} បើបន្ទាត់ $AB$ កែងនឹងប្លង់ $P$ ត្រង់ $B$ ហើយ $BO$ កែងនឹងបន្ទាត់ $a$ នៃប្លង់ $P$ ត្រង់ $O$ នោះបន្ទាត់ $AO$ កែងនឹងបន្ទាត់ $a$ ។ បើបន្ទាត់ $AB$ កែងនឹងប្លង់ $P$ ត្រង់ $B$ ហើយ $AO$ កែងនឹងបន្ទាត់ $a$ នៃប្លង់ $P$ ត្រង់ $O$ នោះបន្ទាត់ $BO$ កែងនឹងបន្ទាត់ $a$ ។
  \item \textbf{ប្លង់មេដ្យាទ័រ ៖} ប្លង់មេដ្យាទ័រនៃអង្កត់មួយគឺជាប្លង់ដែលកែងនឹងអង្កត់នោះ ត្រង់ចំណុចកណ្ដាល ។ គេមានចំណុច $A$, $B$ ស្ថិតនៅសងខាងនៃប្លង់ $P$ ចំណុចផ្សេងគ្នា $C$, $D$, $E$ ស្ថិតនៅក្នុងប្លង់ $P$ ដែល $AC = BC$, $AD = BD$, $AE = BE$ $\Leftrightarrow$ $P$ ប្លង់មេដ្យាទ័រនៃអង្កត់ $AB$ ។
  \item \textbf{ប្លង់កែងនិងចំណោលកែង ៖} គេថាប្លង់ $P$ កែងនឹងប្លង់ $Q$ កាលណាប្លង់ $P$ មានបន្ទាត់មួយកែងនឹងប្លង់ $Q$ ឬប្លង់ $Q$ មានបន្ទាត់មួយកែងនឹងប្លង់ $P$ ។ ចំណោលកែងនៃចំណុចមួយមកលើប្លង់មួយ គឺជាចំណុចប្រសព្វរវាងបន្ទាត់កែងដែលគូសចេញពីចំណុចនោះនឹងប្លង់ ។
\end{enumerate}
\end{note}

\section*{លំហាត់ជំពូកទី ៨}
\begin{enumerate}[1]
  \item $OABCD$ ជាពីរ៉ាមីតកំពូល $O$ (បាត $ABCD$) ។ បន្ទាត់ដែលជាប្រសព្វរវាងប្លង់ $OAB$ និង $OCD$ កាត់ប្លង់ $ABCD$ ត្រង់ចំណុច $H$ ។ ចូរសង់រូបដោយបំពេញបន្ថែមតាមសម្មតិកម្មដែលបានប្រាប់ ។
  \item គេឱ្យចតុមុខ $ABCD$ មួយ ។ $M$, $N$ និង $P$ ជាចំណុចកណ្ដាលរៀងគ្នានៃជ្រុង $BD$, $DC$ និង $CA$ ។ កំណត់ប្រសព្វរវាងប្លង់ $MNP$ និង $DBC$, $DCA$, $ABD$ ។
  \item គេឱ្យប្លង់ $P$ ស្របនឹងប្លង់ $Q$, $AD \parallel BE$, $CF \parallel BE$ ដែល $A, B, C \in P$ និង $D, E, F \in Q$ ។ បង្ហាញថា $\angle BAC = \angle EDF$ ។
  \item គេមានចតុមុខ $ABCD$ មួយ ។ $I$ និង $J$ ជាចំណុចកណ្ដាលរៀងគ្នានៃទ្រនុង $AB$ និង $AC$, $K$ ស្ថិតនៅលើ $AD$ ដែល $AK = \dfrac{2}{3}AD$ ។ បន្ទាត់ $JK$ និង $CD$ ប្រសព្វគ្នាត្រង់ $E$ បន្ទាត់ $IK$ និង $BD$ ប្រសព្វគ្នាត្រង់ $F$ ។ បង្ហាញថា $EF \parallel IJ$ ។
  \item បើ $AB \perp BD$, $\angle ABD = \dfrac{2}{3}x + 56$, $\angle ABC = 2x - 10$ ។ តើ $AB$ កែងនឹងប្លង់ $BCD$ ឬទេ?
  \item គណនាមាឌនៃសូលីតដែលបានមកពីការបង្វិលជុំវិញបន្ទាត់ $AB$ ក្នុងរូបនីមួយៗខាងក្រោម ៖ \quad ក. ចតុកោណកែងដែលមានជ្រុង $AB$, ទទឹង $3$ cm និងអង្កត់ទ្រូង $5$ cm \quad ខ. ត្រីកោណកែងត្រង់ $B$ ដែល $\angle A = 30^\circ$ និងអ៊ីប៉ូតេនុស $4$ cm \quad គ. ការេ $ADBC$ ជ្រុង $5$ cm ដែល $AB$ ជាអង្កត់ទ្រូង ។
  \item $OABC$ ជាតេត្រាអែតដែលមាន $\angle AOB$, $\angle BOC$ និង $\angle COA$ ជាមុំកែង ។ ក្នុងប្លង់ $OBC$ យក $I$ ជាចំណោលកែងនៃ $O$ លើ $BC$ ។ ក្នុងប្លង់ $OAI$ យក $H$ ជាចំណោលកែងនៃ $O$ លើ $AI$ ។ បង្ហាញថា ៖ ក. $AI$ ជាកម្ពស់មួយរបស់ត្រីកោណ $ABC$ \quad ខ. $OH$ កែងនឹងប្លង់ $ABC$ ។
  \item គេមានបន្ទាត់បី $d_1, d_2, d_3$ ដែលពីរៗមិនស្ថិតក្នុងប្លង់តែមួយ ។ បង្ហាញថាមានបន្ទាត់មួយកាត់តាមបន្ទាត់ទាំងពីរនេះ ។
  \item គេឱ្យបន្ទាត់ $d$ ជាបន្ទាត់ប្រសព្វរវាងប្លង់ $P$ និងប្លង់ $Q$ ។ យក $A$ ជាចំណុចនៅក្នុងប្លង់ $P$ ។ $B$ និង $C$ ពីរចំណុចផ្សេងគ្នានៅក្នុងប្លង់ $Q$ ។ ឧបមាថាចំណុច $A$, $B$ និង $C$ មិនស្ថិតនៅលើបន្ទាត់ $d$ ។ បង្ហាញថាចំណុច $A$, $B$ និង $C$ មិនរត់ត្រង់គ្នា ។
  \item គេមានតេត្រាអែត $ABCD$ ហើយ $I, J, K, L$ ជាចំណុចដែល $AI = \dfrac{1}{3}AC$, $AJ = \dfrac{1}{3}AD$, $BK = \dfrac{2}{3}BC$, $BL = \dfrac{2}{3}BD$ ។ \quad ក. បង្ហាញថា ចតុកោណ $IJLK$ ជាចតុកោណព្នាយ \quad ខ. ចំណុច $O$ ជាចំណុចប្រសព្វរវាងបន្ទាត់ $KI$ និង $LJ$ ។ បង្ហាញថាចំណុច $A$, $B$ និង $O$ ស្ថិតនៅលើបន្ទាត់តែមួយ ។
  \item គេឱ្យបន្ទាត់ពីរ $AB$ និង $CD$ មិនស្ថិតនៅក្នុងប្លង់តែមួយ និងចំណុច $M$ និង $N$ ស្ថិតនៅលើបន្ទាត់ $AB$ និង $CD$ រៀងគ្នា ដែល $\dfrac{MA}{MB} = \dfrac{NC}{ND}$ ។ បង្ហាញថា $MN$ ស្របនឹងប្លង់មួយ ។
  \item គេឱ្យបន្ទាត់ $AB$ នៃប្រលេឡូក្រាម $ABCD$ មួយស្ថិតនៅក្នុងប្លង់ $P$ ។ តាមកំពូល $D$ គេគូសបន្ទាត់ $d$ ស្របនឹងប្លង់ $P$ ។ បង្ហាញថា ៖ ក. ប្លង់ $P$ ស្របនឹងប្លង់ដែលកំណត់ដោយបន្ទាត់ $d$ និងចំណុច $C$ \quad ខ. $d$ ស្របនឹងបន្ទាត់ប្រសព្វនៃប្លង់ $P$ និងប្លង់ដែលកំណត់ដោយបន្ទាត់ $d$ និងចំណុច $B$ ។
  \item គេឱ្យ $ABCD$ ជាប្រលេឡូក្រាមនៅក្នុងប្លង់ $P$ មួយ ។ $O$ ជាចំណុចស្ថិតនៅក្រៅប្លង់ $P$ ។ \quad ក. កំណត់ប្រសព្វរវាងប្លង់ $OAD$ និង $OBC$, $OAC$ និង $OBD$ ។ \quad ខ. ចំណុច $A'$ ទៅចំកណ្ដាលជ្រុង $OA$ កំណត់ប្រសព្វរវាងប្លង់ $A'BC$ និង $OAD$, $A'BC$ និង $ODC$ ។
  \item គេឱ្យតេត្រាអែត $ABCD$ យក $I, J, K, L, M$ និង $N$ ជាចំណុចកណ្ដាលរៀងគ្នានៃជ្រុង $AC$, $AD$, $BC$, $BD$, $AB$ និង $CD$ ។ \quad ក. បង្ហាញថា $IJLK$ និង $MJNK$ ជាប្រលេឡូក្រាម ។ ទាញបង្ហាញថា អង្កត់ទ្រូងនៃប្រលេឡូក្រាមទាំងពីរនេះកាត់ត្រង់ចំណុច $G$ មួយ ។ \quad ខ. $A_1$ ជាទីប្រជុំទម្ងន់នៃត្រីកោណ $BCD$ ។ បង្ហាញថា $ABN$, $ACL$ និង $ADK$ កាត់គ្នាតាមបន្ទាត់ $AA_1$ ដែលកាត់តាម $G$ ។
  \item គេមានគូបមួយ ដែលចំណុច $I, J, K$ នៅលើទ្រនុងបីដែលចេញពីកំពូល $A$ ដោយ $AI = AJ = AK = x$ ។ \quad ក. ប្រាប់ឈ្មោះត្រីកោណ $IJK$ \quad ខ. គណនាមាឌនៃ $AIJK$ ជាអនុគមន៍នៃ $x$ ។
  \item នៅក្នុងប្លង់ $P$ មួយមានបន្ទាត់ស្របពីរ $d$ និង $d'$ ។ $A$ ជាចំណុចស្ថិតនៅក្រៅប្លង់ $P$ ដែលចំណោលកែងរបស់វាមកលើបន្ទាត់ $d'$ បានចំណុច $H'$ និងមកលើ $d$ បានចំណុច $H$ ។ \quad ក. បង្ហាញថា $d$ អរត្តូកូណាល់ទៅនឹងប្លង់ $AHH'$ \quad ខ. បង្ហាញថា $HH'$ កែងនឹង $d$ \quad គ. បើ $K$ ជាចំណោលនៃ $A$ មកលើបន្ទាត់ $HH'$ បង្ហាញថា $AK$ កែងនឹងប្លង់ $P$ ។
  \item គេមានកែវមួយមានរាងកោន ។ គេចាក់បារតដែលមានដង់ស៊ីតេ $13.59$ ទឹកមានដង់ស៊ីតេ $1$ និងប្រេងមានដង់ស៊ីតេ $0.915$ ជាបន្តបន្ទាប់គ្នា ។ អង្គធាតុរាវទាំងបីចាក់បំពេញកែវដោយមិនចូលគ្នា ដែលបង្កើតបាននូវស្រទាប់បីផ្សេងគ្នា ហើយមានកម្រាស់ស្មើៗគ្នា ។ តើម៉ាសនៃអង្គធាតុមួយណាធំជាងគេ ។
  \item រកមុំដែលបង្កើតឡើងដោយប្លង់ពីរនៃគូប $ABCDEFGH$ ដូចខាងក្រោម ៖ \quad ក. ប្លង់ $AEGC$ និងប្លង់ $AEHD$ \quad ខ. ប្លង់ $ABCD$ និងប្លង់ $BFHD$ ។
  \item តើសំណើខាងក្រោមមួយណាពិត មួយណាមិនពិត ៖
  \begin{tasks}(2)
    \task $l \parallel \alpha$, $m \parallel \alpha \Rightarrow l \parallel m$
    \task $l \perp \alpha$, $m \perp \alpha \Rightarrow l \parallel m$
    \task $l \perp \alpha$, $m \parallel \alpha \Rightarrow l \perp m$
    \task $\alpha \perp \gamma$, $\beta \perp \gamma \Rightarrow \alpha \parallel \beta$
    \task $\alpha \parallel \beta$, $\beta \parallel \gamma \Rightarrow \alpha \parallel \gamma$
  \end{tasks}
  \item កំណត់ថាតើគែមរបស់តេត្រាអែត $ABCD$ មួយណាដែលបត់ត្រួតស៊ីគ្នាទៅវិញទៅមក ។
  \item គេឱ្យបួនចំណុច $A, B, C, D$ ។ បើបន្ទាត់ $AB$ និង $CD$ មិនត្រួតស៊ីគ្នា នោះបន្ទាត់ $AD$ និង $BC$ មិនត្រួតស៊ីគ្នាដែរ ។ ចូរបង្ហាញ ។
  \item បើប្លង់ពីរ $\alpha$ និង $\beta$ ប្រសព្វគ្នាបានបន្ទាត់ $L$ មួយ ហើយគេគូសបន្ទាត់ $PA$ និង $PB$ ពីរចំណុច $P$ កែងទៅនឹងប្លង់ទាំងពីរខាងលើ ។ ចូរបង្ហាញថា បើបន្ទាត់ $AA'$ និង $BB'$ កែងនឹងបន្ទាត់ $L$ នោះ $A'$ និង $B'$ ត្រួតលើ ។
  \item បង្ហាញថា បើបន្ទាត់ $PA$ មិននៅក្នុងប្លង់ $\alpha$ ទៅបន្ទាត់ $a$ នៅក្នុងប្លង់ $\alpha$ ហើយបន្ទាត់ $PQ$ កែងទៅនឹងបន្ទាត់ $\alpha$ នោះបន្ទាត់ $AQ$ កែងទៅនឹងប្លង់ $\alpha$ ។
  \item បង្ហាញថា បើបន្ទាត់ $PQ$ មិននៅក្នុងប្លង់ $\alpha$ ទៅប្លង់ $\alpha$ នៅក្នុងប្លង់ $\alpha$ ហើយបន្ទាត់ $PQ$ កែងទៅនឹងបន្ទាត់ $a$ នោះបន្ទាត់ $PA$ កែងទៅនឹងបន្ទាត់ $a$ ។
  \item បង្ហាញថា បើបន្ទាត់ $L$ កែងទៅនឹងបន្ទាត់ពីរប្រសព្វគ្នានៅក្នុងប្លង់ $P$ នោះបន្ទាត់ $L$ កែងទៅនឹងប្លង់នោះ ។
  \item រកមុំដែលកើតឡើងដោយជ្រុងពីរនៃគូប $ABCDEFGH$ ដូចខាងក្រោម ៖ \quad ក. $BC$ និង $GH$ \quad ខ. $BA$ និង $FC$ \quad គ. $BD$ និង $AH$ ។
  \item តើអ្វីជាចំណុចប្រសព្វនៃប្លង់ $MAB$, $MBC$ និង $MAC$?
  \item គូសបន្ទាត់ $PA$ ពីចំណុច $P$ ដែលមិននៅក្នុងប្លង់ $M$ ហើយកែងទៅនឹងបន្ទាត់ $a$ នៅក្នុងប្លង់ $M$ ។ នៅក្នុងស្ថានភាពនេះ បើយើងគូសបន្ទាត់ $PQ$ ពីចំណុច $P$ កែងទៅនឹងបន្ទាត់ $b$ នោះបន្ទាត់ $PQ$ កែងទៅនឹងប្លង់ $M$ ។
  \item បង្ហាញថា បើបួនចំណុច $A, B, C, D$ នៅក្នុងលំហដែលផ្ទៀងផ្ទាត់ $AB^2 + CD^2 = AC^2 + BD^2$ នោះ $AD$ កែងនឹង $BC$ ។
  \item គេឱ្យត្រីកោណ $ABC$ មួយ ហើយ $AD$ ជាកន្លះបន្ទាត់ពុះមុំ $BAC$ ។ បង្ហាញថា $AB : AC = BD : CD$ ។
  \item គេឱ្យត្រីកោណ $ABC$ មួយ ហើយ $AD$ ជាមេដ្យាន ដែល $AD = BD = CD$ ។ បង្ហាញថាត្រីកោណ $ABC$ ជាត្រីកោណកែង ។
  \item រកចំនួនបន្ទាត់ដែលកាត់តាមចំណុច $B$ នៃពីរ៉ាមីត $MABCD$ (បន្ទាត់កាត់តាមកំពូលពីរនៃពីរ៉ាមីត) ។
  \item រកប្លង់ចំនួនបីដែលកាត់តាមចំណុចទាំងពីរ $A$ និង $B$ នៃប្រលេពីប៉ែតកែង $ABCDA'B'C'D'$ ។
  \item ពិចារណាអំពីព្រីសរាងចតុកោណ $ABCDEFGH$ (មុខលើ $ABCD$ និងមុខក្រោម $EFGH$ ដែល $E, F, G, H$ នៅក្រោម $A, B, C, D$ រៀងគ្នា) ដែល $AB \ne BC \ne AE$ ។
  \begin{enumerate}[a]
    \item តើអាចនិយាយអ្វីខ្លះអំពី $EF$ និង $FG$? (ក. វាស្របគ្នា \quad ខ. វាកាត់កែង ហើយកែងគ្នា \quad គ. ពួកវាមិនស្របឬកាត់ហើយមិនកែងគ្នា \quad ឃ. វាមិនកាត់គ្នានិងកែងគ្នាទេ)
    \item តើអាចនិយាយអ្វីខ្លះអំពី $AE$ និង $CG$? (ក. វាមិនកាត់គ្នានិងកែងគ្នាទេ \quad ខ. វាកាត់កែងនិងកែងគ្នា \quad គ. ពួកវាស្របគ្នា \quad ឃ. វាមិនស្របឬមិនកាត់គ្នានិងកែងគ្នាទេ)
    \item តើអាចនិយាយអ្វីខ្លះអំពី $BD$ និង $DH$? (ក. វាកាត់គ្នានិងកែងគ្នា \quad ខ. ពួកវាមិនស្របឬកាត់កែងគ្នា \quad គ. វាមិនកាត់គ្នានិងកែងគ្នាទេ \quad ឃ. វាស្របគ្នា)
    \item តើអាចនិយាយអ្វីខ្លះអំពី $BD$ និង $AC$? (ក. វាស្របគ្នា \quad ខ. វាកាត់គ្នានិងកែងគ្នា \quad គ. វាមិនស្របគ្នាឬកាត់គ្នា និងកែងគ្នាទេ \quad ឃ. វាមិនកាត់គ្នានិងកែងគ្នាទេ)
  \end{enumerate}
  \item អ្នកមានប្រអប់ដែលមានទំហំ $2$ ហ្វីត $1$ ហ្វីត ទទឹង $6$ អ៊ីង ។ តើផលិតផលប៉ុន្មានដែលសមនឹងខាងក្នុង ប្រសិនបើកញ្ចប់ផលិតផលនីមួយៗវាស់ $3$ អ៊ីញ គុណនឹង $6$ អ៊ីង និងកម្ពស់ $1$ អ៊ីញ?
  \item គេឱ្យបន្ទាត់ $L$ មួយកាត់ប្លង់ $X$ ត្រង់ចំណុច $A$ ។ ចូរពិនិត្យមើលប្រយោគខាងក្រោមថា តើមួយណាត្រឹមត្រូវ? \quad ក. បន្ទាត់ស្របទៅនឹងប្លង់ \quad ខ. បន្ទាត់ស្ថិតនៅក្នុងប្លង់ \quad គ. បន្ទាត់កាត់ប្លង់ ។
\end{enumerate}

\section*{លំហាត់អនុវត្តន៍}
\begin{enumerate}[1]
  \item បញ្ជាក់ ត្រូវ ឬខុស ៖ \quad ក. នៅក្នុងលំហ បើគេមានបន្ទាត់ពីរ $a$ និង $a'$ អរត្តូកូណាល់ទៅនឹងបន្ទាត់ $b$ នោះគេថា $a \parallel a'$ ។ \quad ខ. បន្ទាត់ $d$ និង $d'$ អរត្តូកូណាល់ទៅនឹងប្លង់ $P$ មួយ នោះគេបាន $d \parallel d'$ ។ \quad គ. គេមានបន្ទាត់ពីរ $a$ និង $a'$ ។ បើ $a$ អរត្តូកូណាល់ទៅនឹងប្លង់ $P$ និង $a'$ ស្របទៅនឹងប្លង់ $P$ នោះ $a$ អរត្តូកូណាល់ទៅនឹង $a'$ ។ \quad ឃ. បន្ទាត់មួយអរត្តូកូណាល់ទៅនឹងបន្ទាត់ពីរផ្សេងគ្នា ដែលស្ថិតនៅក្នុងប្លង់ $P$ មួយ នោះវាអរត្តូកូណាល់ទៅនឹងប្លង់ ។ \quad ង. បើគេមានបន្ទាត់ពីរ $d$ និង $d'$ មិនស្របគ្នា ហើយអរត្តូកូណាល់រៀងគ្នាទៅនឹងប្លង់ $P$ និង $P'$ នោះគេថាប្លង់ $P$ និង $P'$ ប្រសព្វគ្នា ។
  \ans
  ក. ខុស \quad ខ. ត្រូវ \quad គ. ត្រូវ \quad ឃ. ខុស (បន្ទាត់ទាំងពីរត្រូវតែប្រសព្វគ្នា) \quad ង. ត្រូវ ។
  \item គេមានគូប $ABCDA'B'C'D'$ មួយដែល $M$, $N$, $P$, $Q$ ជាចំណុចកណ្ដាលនៃទ្រនុង $AB$, $CD$, $A'D'$, $C'D'$ រៀងគ្នា ។ ចូរបញ្ជាក់ ត្រូវ ឬខុស ៖ \quad ក. បន្ទាត់ $MN \parallel PN$ \quad ខ. បន្ទាត់ $MN$ អរត្តូកូណាល់នឹង $PN$ \quad គ. បន្ទាត់ $MN$ មិនកាត់និងមិនស្របនឹង $PN$ \quad ឃ. ប្លង់មេដ្យាទ័រនៃអង្កត់ $MN$ គឺជាប្លង់ $BDA'$ ។
  \ans
  ក. ខុស \quad ខ. ខុស \quad គ. ខុស ($MN$ និង $PN$ ប្រសព្វគ្នាត្រង់ $N$) \quad ឃ. ខុស (ប្លង់មេដ្យាទ័រនៃ $MN$ កាត់តាមចំណុចកណ្ដាលនៃ $AD$, $BC$, $A'D'$, $B'C'$ ហើយស្របនឹងមុខ $ABB'A'$) ។ (សម្គាល់ ៖ បន្ទាត់ $MN$ និង $PQ$ មិនកាត់និងមិនស្របគ្នា ។)
  \item $SABCD$ ជាពីរ៉ាមីតនិយ័តមួយដែលមានបាតជាការេ ។ ចូរបំពេញក្នុងតារាងខាងក្រោម ៖
  \[
  \begin{array}{|l|c|c|c|c|}\hline
  \text{កម្ពស់} & 1.5\ \text{m} & & 0.6\ \text{m} & \\ \hline
  \text{ជ្រុងការេ} & 0.5\ \text{m} & & & 3\ \text{m}\\ \hline
  \text{ផ្ទៃក្រឡាបាត} & & 144\ \text{dm}^2 & & \\ \hline
  \text{មាឌ} & & 0.624\ \text{dm}^3 & 242\ \text{m}^3 & 36\ \text{m}^3\\ \hline
  \end{array}
  \]
  \ans
  ប្រើ $V = \dfrac{1}{3}Bh$ ៖ ជួរទី ១ ៖ $B = 0.25\ \text{m}^2$, $V = 0.125\ \text{m}^3$ ។ ជួរទី ២ ៖ ជ្រុង $12$ dm, $h = \dfrac{3\times 0.624}{144} = 0.013$ dm ។ ជួរទី ៣ ៖ $B = \dfrac{3\times 242}{0.6} = 1210\ \text{m}^2$, ជ្រុង $= 11\sqrt{10} \approx 34.79$ m ។ ជួរទី ៤ ៖ $B = 9\ \text{m}^2$, $h = \dfrac{3\times 36}{9} = 12$ m ។
  \item គេមានកោណមួយនិងប្លង់មួយកាត់តាមចំណុចកណ្ដាលនៃកម្ពស់ រួចស្របទៅនឹងប្លង់បាត ។ ចូរបំពេញចម្លើយក្នុងតារាងខាងក្រោម ៖
  \[
  \begin{array}{|l|c|c|c|c|c|}\hline
   & \text{កម្ពស់} & \text{កាំនៃបាត} & \text{ជនេត្រ} & \text{ផ្ទៃក្រឡាបាត} & \text{មាឌ}\\ \hline
  \text{កោណមិនទាន់កាត់} & h & & a & \pi r^2 & v\\ \hline
  \text{កោណបានកាត់រួច} & & \dfrac{r}{2} & & & \\ \hline
  \end{array}
  \]
  \ans
  កោណមិនទាន់កាត់ ៖ កាំបាត $r$ ។ កោណតូចដែលបានកាត់រួច ៖ កម្ពស់ $\dfrac{h}{2}$, កាំ $\dfrac{r}{2}$, ជនេត្រ $\dfrac{a}{2}$, ផ្ទៃក្រឡាបាត $\dfrac{\pi r^2}{4}$, មាឌ $\dfrac{v}{8}$ ។
  \item គេមានចតុមុខនិយ័ត $ABCD$ មួយ ។ បង្ហាញថាជ្រុង $CD$ ស្ថិតនៅក្នុងប្លង់មេដ្យាទ័រនៃអង្កត់ $AB$ ។
  \ans
  $CA = CB$ និង $DA = DB$ (គ្រប់ទ្រនុងស្មើគ្នា) ដូច្នេះ $C$ និង $D$ ស្ថិតក្នុងប្លង់មេដ្យាទ័រនៃ $AB$ ។ ហេតុនេះបន្ទាត់ $CD$ ស្ថិតក្នុងប្លង់នោះ ។
  \item បន្ទាត់ $d$ និង $\ell$ មិនស្ថិតនៅក្នុងប្លង់តែមួយ ប៉ុន្តែអរត្តូកូណាល់នឹងគ្នា ។ ចំណុច $A$ និង $H$ ស្ថិតនៅលើ $\ell$ និង $d$ រៀងគ្នាដែល $H$ ជាចំណោលកែងនៃ $A$ មកលើ $d$ ។ បង្ហាញថា $d$ កែងនឹងប្លង់ដែលកំណត់ដោយ $H$ និង $\ell$ ។
  \ans
  $d \perp \ell$ (សម្មតិកម្ម) និង $d \perp AH$ ។ បន្ទាត់ $\ell$ និង $AH$ ប្រសព្វគ្នាត្រង់ $A$ ហើយស្ថិតក្នុងប្លង់ដែលកំណត់ដោយ $H$ និង $\ell$ ។ ដូច្នេះ $d$ កែងនឹងប្លង់នោះ ។
  \item គេឱ្យពីរ៉ាមីត $SABCD$ ដែល $ABCD$ ជាការេ និង $SB$ កែងនឹងប្លង់ $ABCD$ ។ \quad ក. បង្ហាញថាបន្ទាត់ $AD$ កែងនឹងប្លង់ $ABS$ \quad ខ. បើ $AB = BS = a$ គណនា $BD$, $DS$, $SA$ ជាអនុគមន៍នៃ $a$ ។ រកមាឌនៃពីរ៉ាមីត $SABCD$ ជាអនុគមន៍នៃ $a$ ។
  \ans
  ក. $AD \perp AB$ និង $AD \perp SB$ ($SB \perp ABCD$) ដូច្នេះ $AD \perp (ABS)$ ។ \quad ខ. $BD = a\sqrt{2}$, $DS = \sqrt{SB^2 + BD^2} = a\sqrt{3}$, $SA = a\sqrt{2}$ ។ $V = \dfrac{1}{3}a^2\cdot a = \dfrac{a^3}{3}$ ។
  \item គេឱ្យបន្ទាត់ $d$ ជាប្រសព្វរវាងប្លង់ $P$ និង $P'$ និងចំណុច $A$ មួយមិនស្ថិតនៅក្នុងប្លង់ទាំងពីរ ។ បន្ទាត់ពីរប្រសព្វគ្នាត្រង់ចំណុច $A$ កែងនឹងប្លង់ $P$ និង $P'$ ត្រង់ $H$ និង $H'$ រៀងគ្នា ។ \quad ក. បង្ហាញថា $d$ អរត្តូកូណាល់នឹងប្លង់ $AHH'$ \quad ខ. តាងចំណុច $K$ ជាប្រសព្វរវាងបន្ទាត់ $d$ និងប្លង់ $AHH'$ ។ បង្ហាញថា $AK$ កែងនឹងបន្ទាត់ $d$ \quad គ. ឧបមាថា $AH = AH'$ បង្ហាញថា $KH = KH'$ ។
  \ans
  ក. $d \subset P$ នាំឱ្យ $d \perp AH$ ។ $d \subset P'$ នាំឱ្យ $d \perp AH'$ ។ ដូច្នេះ $d \perp (AHH')$ ។ \quad ខ. $AK \subset (AHH')$ នាំឱ្យ $AK \perp d$ ។ \quad គ. ត្រីកោណកែង $AHK$ និង $AH'K$ (កែងត្រង់ $H$ និង $H'$) មានអ៊ីប៉ូតេនុសរួម $AK$ និង $AH = AH'$ ដូច្នេះ $KH = KH'$ ។
  \item គេមានត្រីកោណកែង $ABC$ កែងត្រង់ $A$ និង $AB = a$, $AC = 2a$ ។ យក $d$ ជាបន្ទាត់មួយកែងនឹងប្លង់ $ABC$ ត្រង់ $A$ ។ ជ្រើសរើសចំណុច $D$ មួយនៅលើ $d$ ដែល $AD = 2a$ ។ \quad ក. បង្ហាញថា $AC$ អរត្តូកូណាល់នឹង $BD$ និង $AB$ អរត្តូកូណាល់នឹង $CD$ \quad ខ. $I$ ជាចំណុចកណ្ដាលនៃអង្កត់ $CD$ ។ បង្ហាញថាប្លង់ $ABI$ ជាប្លង់មេដ្យាទ័រនៃអង្កត់ $CD$ \quad គ. គណនាមាឌ $ABCD$ ជាអនុគមន៍នៃ $a$ ។
  \ans
  ក. $AC \perp AB$ និង $AC \perp AD$ នាំឱ្យ $AC \perp (ABD)$ ដូច្នេះ $AC \perp BD$ ។ ដូចគ្នា $AB \perp (ACD)$ ដូច្នេះ $AB \perp CD$ ។ \quad ខ. $AC = AD = 2a$ និង $BC = BD = a\sqrt{5}$ ហើយ $IC = ID$ ។ ចំណុច $A, B, I$ មិនរត់ត្រង់គ្នា ហើយស្មើចម្ងាយពី $C$ និង $D$ ដូច្នេះ $ABI$ ជាប្លង់មេដ្យាទ័រនៃ $CD$ ។ \quad គ. $V = \dfrac{1}{3}\times\dfrac{1}{2}\times a\times 2a\times 2a = \dfrac{2a^3}{3}$ ។
  \item រង្វង់ផ្ចិត $O$ មួយស្ថិតនៅក្នុងប្លង់ $P$ ដែលមានអង្កត់ផ្ចិត $AB$ ។ $C$ ជាចំណុចមួយនៅលើរង្វង់ផ្សេងពីចំណុច $A$ និង $B$ និង $d$ ជាបន្ទាត់កាត់តាម $A$ ហើយអរត្តូកូណាល់នឹងប្លង់ $P$ ។ យកចំណុច $D$ ឋិតនៅលើបន្ទាត់ $d$ ផ្សេងពី $A$ ។ \quad ក. បង្ហាញថា $BC$ កែងនឹងប្លង់ $ACD$ \quad ខ. $H$ ជាចំណុចមួយនៅលើ $CD$ ដែល $AH$ កែងនឹង $CD$ ។ បង្ហាញថា $AH$ កែងនឹង $BCD$ \quad គ. បង្ហាញថា $AB^2 = AH^2 + HC^2 + CB^2$ ។
  \ans
  ក. $BC \perp AC$ (មុំចារឹកក្នុងកន្លះរង្វង់) និង $BC \perp AD$ ($AD \perp P$) ដូច្នេះ $BC \perp (ACD)$ ។ \quad ខ. $AH \subset (ACD)$ នាំឱ្យ $AH \perp BC$ ហើយ $AH \perp CD$ ដូច្នេះ $AH \perp (BCD)$ ។ \quad គ. $AH \perp HB$ នាំឱ្យ $AB^2 = AH^2 + HB^2$ ។ $BC \perp (ACD)$ នាំឱ្យ $BC \perp HC$ ហើយ $HB^2 = HC^2 + CB^2$ ។ ដូច្នេះ $AB^2 = AH^2 + HC^2 + CB^2$ ។
\end{enumerate}
